% ─────────────────────────────────────────────────────────────────────────────
%  LAB 02 CHEATSHEET — Arithmetic & Logical Ops + 2D Convolution
%  Bùi Quang Chiến — 23001837   |   MAT3562E Computer Vision 2025-2026
% ─────────────────────────────────────────────────────────────────────────────
\documentclass[9pt,a4paper,landscape]{extarticle}

\usepackage[utf8]{inputenc}
\usepackage[T1]{fontenc}
\usepackage[margin=1cm]{geometry}
\usepackage{amsmath,amssymb}
\usepackage{multicol}
\usepackage{booktabs}
\usepackage{xcolor}
\usepackage{listings}
\usepackage{mdframed}
\usepackage{titlesec}
\usepackage{enumitem}
\usepackage{parskip}

\setlength{\parindent}{0pt}
\setlength{\columnsep}{1em}
\setlength{\columnseprule}{0.3pt}
\setlist{nosep, leftmargin=1em}

% ─── Colors ──────────────────────────────────────────────────────────────────
\definecolor{hdr}{RGB}{30,80,150}
\definecolor{subhdr}{RGB}{60,120,200}
\definecolor{codebg}{RGB}{245,245,248}
\definecolor{boxbg}{RGB}{230,242,255}
\definecolor{warnbg}{RGB}{255,243,220}
\definecolor{greenbg}{RGB}{230,255,235}


% ─── Code style ──────────────────────────────────────────────────────────────
\lstset{
  basicstyle=\ttfamily\scriptsize,
  backgroundcolor=\color{codebg},
  frame=none,
  breaklines=true,
  tabsize=2,
  showstringspaces=false,
  keywordstyle=\color{blue!70}\bfseries,
  commentstyle=\color{green!50!black},
  stringstyle=\color{red!60!black},
  aboveskip=1pt, belowskip=1pt,
}

% ─── Note boxes ──────────────────────────────────────────────────────────────
\newmdenv[backgroundcolor=warnbg, linewidth=0.5pt, linecolor=orange!60, innerleftmargin=4pt, innerrightmargin=4pt, innertopmargin=2pt, innerbottommargin=2pt, skipabove=1pt, skipbelow=1pt]{notebox}
\newmdenv[backgroundcolor=greenbg, linewidth=0.5pt, linecolor=green!50!black, innerleftmargin=4pt, innerrightmargin=4pt, innertopmargin=2pt, innerbottommargin=2pt, skipabove=1pt, skipbelow=1pt]{tipbox}
\newmdenv[backgroundcolor=boxbg, linewidth=0.5pt, linecolor=hdr, innerleftmargin=4pt, innerrightmargin=4pt, innertopmargin=2pt, innerbottommargin=2pt, skipabove=1pt, skipbelow=1pt]{infobox}

% ─────────────────────────────────────────────────────────────────────────────
\begin{document}
\begin{center}
  {\large\bfseries\color{hdr} LAB 02 CHEATSHEET — Arithmetic \& Logical Ops + 2D Convolution} \\
  {\small Bui Quang Chien \textbullet{} 23001837 \textbullet{} MAT3562E Computer Vision 2025--2026}
\end{center}
\hrule\vspace{2pt}

\begin{multicols}{3}

% ══════════════════════════════════════════════════════════════════════════════
\section{Core Utilities}

\subsection{Type conversion (ALWAYS use these)}
\begin{lstlisting}[language=Python]
# uint8 [0,255] -> float32 [0,1]
r = img_u8.astype(np.float32) / 255.0

# float [0,1] -> uint8 [0,255]  (clip first!)
out = (np.clip(f, 0, 1) * 255 + 0.5).astype(np.uint8)

# any float -> uint8 (min-max normalise)
def normalise(arr):
    mn, mx = arr.min(), arr.max()
    if mx == mn: return np.zeros_like(arr, np.uint8)
    return ((arr-mn)/(mx-mn)*255).astype(np.uint8)
\end{lstlisting}

\begin{notebox}
\textbf{Safe workflow:} uint8 $\to$ float $\to$ compute $\to$ clip $\to$ uint8. \\
Direct uint8 arithmetic with NumPy \textbf{wraps around} (overflow/underflow)!
\end{notebox}

\subsection{Read / Show / Save}
\begin{lstlisting}[language=Python]
img = cv2.imread(path, cv2.IMREAD_GRAYSCALE)
plt.imshow(img, cmap='gray', vmin=0, vmax=255)
cv2.imwrite(path, img)
\end{lstlisting}

% ══════════════════════════════════════════════════════════════════════════════
\section{Part 1A — Gamma Correction}

\textbf{Formula:} $s = r^{\gamma},\quad r \in [0,1]$

\begin{center}
\begin{tabular}{@{}cc@{}}
\toprule
$\gamma$ & Effect \\
\midrule
$< 1$ & Brighter (raise shadows) \\
$= 1$ & Identity \\
$> 1$ & Darker (compress shadows) \\
\bottomrule
\end{tabular}
\end{center}

\begin{lstlisting}[language=Python]
# NumPy
s = np.power(to_float01(img), gamma)
result = to_uint8(s)

# OpenCV (LUT = precomputed table, faster)
table = np.array([((i/255)**gamma)*255
                   for i in range(256)], dtype=np.uint8)
result = cv2.LUT(img_u8, table)
\end{lstlisting}

\begin{tipbox}
\textbf{When to use:} Preprocess before Otsu on dark images ($\gamma < 1$).
Can make histogram more bimodal $\Rightarrow$ better Otsu threshold.
\end{tipbox}

% ══════════════════════════════════════════════════════════════════════════════
\section{Part 1B — Arithmetic Operations}

$S_+ = \alpha A + \beta B$, \ $S_- = \alpha A - \beta B$, \
$S_\times = A \odot B$, \ $S_\div = \frac{A}{B+\varepsilon}$

\begin{lstlisting}[language=Python]
# NumPy (float + clip) — flexible
A, B = to_float01(a), to_float01(b)
add_  = to_uint8(A + B)           # clip at 255
sub_  = to_uint8(A - B)           # clip at 0
mul_  = to_uint8(A * B)           # always <= 1
div_  = to_uint8(A / (B + 1e-6))  # can exceed 1!

# OpenCV (saturating uint8) — fast
cv2.add(a, b)          # saturates at 255
cv2.subtract(a, b)     # saturates at 0
cv2.multiply(a, b, scale=...)
cv2.divide(a, b, scale=...)
\end{lstlisting}

\textbf{Key differences:}
\begin{itemize}
  \item NumPy: soft clip (lossy but controlled)
  \item OpenCV: saturating integer (no wrapping, but less flexible)
  \item Division: clipping \textbf{hides} meaningful signal; prefer \texttt{normalise()}
\end{itemize}

\textbf{Change detection:} $|A - B|$ is near-zero everywhere the scene is unchanged.

% ══════════════════════════════════════════════════════════════════════════════
\section{Part 1C — Logical Operations}

\textbf{Best on binary images (0 / 255)}

\begin{lstlisting}[language=Python]
# NumPy
np.where((a==255) & (b==255), 255, 0)  # AND
np.where((a==255) | (b==255), 255, 0)  # OR
np.where((a==255) ^ (b==255), 255, 0)  # XOR

# OpenCV
cv2.bitwise_and(a, b)
cv2.bitwise_or(a, b)
cv2.bitwise_xor(a, b)
\end{lstlisting}

\begin{center}
\begin{tabular}{@{}ll@{}}
\toprule
Goal & Operator \\
\midrule
Consensus (fewest FP) & AND \\
Union (fewest FN) & OR \\
Disagreement / edges & XOR \\
Extract ROI & AND(mask, original) \\
\bottomrule
\end{tabular}
\end{center}

% ══════════════════════════════════════════════════════════════════════════════
\section{Part 1D — Otsu Thresholding}

\textbf{Goal:} Automatically find $t^*$ that maximises between-class variance:
\[
  \sigma_b^2(t) = \omega_0 \omega_1 (\mu_0 - \mu_1)^2
\]

\begin{lstlisting}[language=Python]
# NumPy (from scratch)
hist = np.bincount(img.ravel(), minlength=256).astype(np.float64)
p     = hist / hist.sum()
omega = np.cumsum(p)
mu    = np.cumsum(p * np.arange(256))
mu_T  = mu[-1]
w0, w1 = omega, 1-omega
mu0 = mu/(w0+1e-12);  mu1=(mu_T-mu)/(w1+1e-12)
t_star = int(np.argmax(w0*w1*(mu0-mu1)**2))

# OpenCV
t, bimg = cv2.threshold(img, 0, 255,
               cv2.THRESH_BINARY + cv2.THRESH_OTSU)
\end{lstlisting}

\begin{notebox}
\textbf{Works best:} near-bimodal histogram, uniform lighting, moderate noise. \\
\textbf{Fails:} non-uniform illumination, many histogram modes, tiny foreground.
\end{notebox}

% ══════════════════════════════════════════════════════════════════════════════
\section{Integrated Pipeline (Part 1)}

\begin{lstlisting}[language=Python]
# Exercise 1: Brighten then Otsu
step1 = gamma_correction_np(img, 0.6)
t, mask = cv2.threshold(step1, 0, 255,
              cv2.THRESH_BINARY + cv2.THRESH_OTSU)

# Exercise 2: ROI Extraction
roi = cv2.bitwise_and(img_gray, mask)

# Exercise 3: Change Detection
diff = cv2.absdiff(imgA, imgB)
t, change_mask = cv2.threshold(diff, 0, 255,
                     cv2.THRESH_BINARY + cv2.THRESH_OTSU)

# Exercise 4: Mask Fusion
or_mask  = cv2.bitwise_or(mask1, mask2)   # candidates
and_mask = cv2.bitwise_and(mask1, mask2)  # consensus
xor_mask = cv2.bitwise_xor(mask1, mask2)  # disagreement
\end{lstlisting}

% ══════════════════════════════════════════════════════════════════════════════
\section{Part 2 — 2D Convolution Theory}

\[
  y[i,j] = \sum_{u}\sum_{v} K[u,v]\cdot x[i+u-\Delta_h,\; j+v-\Delta_w]
\]
$\Delta_h = \lfloor k_H / 2 \rfloor$, \ $\Delta_w = \lfloor k_W / 2 \rfloor$

\begin{infobox}
\textbf{Convolution vs Correlation:} True conv flips kernel $180^\circ$.
\texttt{cv2.filter2D} = cross-correlation (no flip).
For \textbf{symmetric kernels}: identical. Flip manually for asymmetric (e.g.\ Sobel).
\end{infobox}

\subsection{Padding Modes}

\begin{center}
\small
\begin{tabular}{@{}lll@{}}
\toprule
Mode & Output size & Border \\
\midrule
VALID & $(H-k+1)\times(W-k+1)$ & Discarded \\
SAME zero & $H\times W$ & Zeros outside \\
SAME reflect & $H\times W$ & Mirrored \\
SAME replicate & $H\times W$ & Repeat edge pixel \\
\bottomrule
\end{tabular}
\end{center}

After a single $3\times3$ VALID conv: size shrinks by 2. \\
After $n$ such layers: shrinks by $2n$. $\Rightarrow$ CNNs use SAME.

% ══════════════════════════════════════════════════════════════════════════════
\section{Part 2 — NumPy Implementation}

\begin{lstlisting}[language=Python]
# VALID (no padding)
def conv2d_valid_np(img_u8, kernel):
    img_f  = img_u8.astype(np.float32)
    k      = kernel[::-1,::-1]       # flip for true conv
    kH,kW  = k.shape
    oH,oW  = img_f.shape[0]-kH+1, img_f.shape[1]-kW+1
    out    = np.zeros((oH,oW), np.float64)
    for i in range(oH):
        for j in range(oW):
            out[i,j] = np.sum(img_f[i:i+kH,j:j+kW]*k)
    return normalise(out)

# SAME (with padding)
def conv2d_same_np(img_u8, kernel, pad='reflect'):
    img_f  = img_u8.astype(np.float32)
    k      = kernel[::-1,::-1]
    kH,kW  = k.shape
    pH,pW  = kH//2, kW//2
    padded = np.pad(img_f,((pH,pH),(pW,pW)), mode=pad)
    H,W    = img_f.shape
    out    = np.zeros((H,W), np.float64)
    for i in range(H):
        for j in range(W):
            out[i,j]=np.sum(padded[i:i+kH,j:j+kW]*k)
    return normalise(out)
\end{lstlisting}

% ══════════════════════════════════════════════════════════════════════════════
\section{Part 2 — OpenCV Implementation}

\begin{lstlisting}[language=Python]
# SAME mode, normalised display
def conv2d_cv2(img_u8, kernel,
               border=cv2.BORDER_REFLECT):
    out = cv2.filter2D(img_u8.astype(np.float32),
                       cv2.CV_32F,
                       kernel.astype(np.float32),
                       borderType=border)
    return normalise(out.astype(np.float64))

# borderType options:
# cv2.BORDER_CONSTANT   (zero pad)
# cv2.BORDER_REFLECT    (mirror)
# cv2.BORDER_REPLICATE  (edge repeat)
\end{lstlisting}

% ══════════════════════════════════════════════════════════════════════════════
\section{Part 2 — Kernel Reference}

\begin{lstlisting}[language=Python]
K_BOX3  = np.ones((3,3))/9             # avg blur 3x3
K_GAUSS3= np.array([[1,2,1],[2,4,2],
                    [1,2,1]])/16.0     # Gaussian
K_SHARPEN=np.array([[0,-1,0],
                    [-1,5,-1],[0,-1,0]]) # sharpen
K_SOBEL_X=np.array([[-1,0,1],
                    [-2,0,2],[-1,0,1]]) # horiz grad
K_SOBEL_Y=np.array([[-1,-2,-1],
                    [0,0,0],[1,2,1]])   # vert grad
K_LAP   = np.array([[0,1,0],[1,-4,1],
                    [0,1,0]])           # Laplacian
\end{lstlisting}

\begin{center}
\small
\begin{tabular}{@{}ll@{}}
\toprule
Kernel type & Frequency response \\
\midrule
Average / Gaussian & Low-pass (blur) \\
Sharpen & High-pass (boost edges) \\
Sobel X/Y & Band-pass directional \\
Laplacian & High-pass isotropic \\
LoG & Gaussian + Laplacian \\
\bottomrule
\end{tabular}
\end{center}

\subsection{Sobel Magnitude}
\begin{lstlisting}[language=Python]
Gx = cv2.filter2D(img.astype(np.float32),
                  cv2.CV_32F, K_SOBEL_X)
Gy = cv2.filter2D(img.astype(np.float32),
                  cv2.CV_32F, K_SOBEL_Y)
mag = np.sqrt(Gx**2 + Gy**2)
mag_u8 = normalise(mag)
\end{lstlisting}

% ══════════════════════════════════════════════════════════════════════════════
\section{Part 2 — Full Pipeline}

\begin{lstlisting}[language=Python]
# Blur -> Gamma -> Otsu -> ROI
step1 = conv2d_same_np(img, K_GAUSS5, 'reflect')
step2 = gamma_correction_np(step1, 0.6)
t, mask = cv2.threshold(step2, 0, 255,
              cv2.THRESH_BINARY+cv2.THRESH_OTSU)
roi = cv2.bitwise_and(img, mask)
\end{lstlisting}

\begin{tipbox}
\textbf{Why blur first?} Removes texture noise $\Rightarrow$ histogram has fewer
spurious peaks $\Rightarrow$ Otsu finds a cleaner boundary.
\end{tipbox}

% ══════════════════════════════════════════════════════════════════════════════
\section{Key Theory to Remember}

\textbf{Gamma correction:}
$s = r^\gamma$. Fixed points at $r=0$ and $r=1$.
$\gamma < 1$ raises midtones; $\gamma > 1$ lowers them.

\textbf{Otsu:} Maximises $\sigma_b^2 = \omega_0\omega_1(\mu_0-\mu_1)^2$.
Assumes bimodal histogram. Equivalent to minimising within-class variance.

\textbf{Convolution = weighted sliding window.} Kernel determines the operation type:
\begin{itemize}
  \item Sum of kernel = 1 $\Rightarrow$ no brightness change (blur)
  \item Sum of kernel = 0 $\Rightarrow$ detects change (edges)
  \item Sum of kernel > 1 $\Rightarrow$ brightens output
\end{itemize}

\textbf{VALID vs SAME:}
VALID output size = $(H - k + 1) \times (W - k + 1)$.
SAME output size = $H \times W$ (padding required).

\textbf{Box vs Gaussian blur:}
Box blur: equal weights → ringing artefacts.
Gaussian: smooth bell-curve weights → fewer artefacts.

\textbf{Laplacian is noisy.}
Always apply Gaussian first (= LoG) for reliable edge detection.

\textbf{Clip vs Normalise for display:}
\begin{itemize}
  \item Clip: correct for filters where output ∈ [0,255] (blur, sharpen)
  \item Normalise: needed for edge filters (contain negatives / values > 255)
\end{itemize}

\textbf{np.convolve / scipy.signal vs cv2.filter2D:}
NumPy/SciPy flipkernels. OpenCV does not. For symmetric kernels: same result.

% ══════════════════════════════════════════════════════════════════════════════
\section{Common Mistakes to Avoid}

\begin{notebox}
\begin{itemize}
  \item Doing \texttt{img1 + img2} on \texttt{uint8} directly wraps around (e.g.\ $200 + 100 = 44$). Always convert to float first.
  \item Using \texttt{VALID} mode then expecting the same spatial size as the input.
  \item Applying Laplacian directly to a noisy image (amplifies noise). Use LoG instead.
  \item Using \texttt{cv2.filter2D} for asymmetric kernels without flipping (cross-correlation ≠ convolution for asymmetric kernels).
  \item Normalising sharpened output: use clip, not normalise, or sharpening appears to change brightness.
\end{itemize}
\end{notebox}

\end{multicols}
\end{document}
